\clearpage
\begingroup
% chapter 4 accepted by the user, 27 Sep 2026
\section{The net of a classical phase space}\label{sec:phasespace}

\noindent\motiv{\vet{The observables of the free field are Weyl operators labelled by the
elements of its classical phase space, and an observable belongs to a region when the source that
generates the phase-space element can be chosen inside that region. This \revh{chapter} states that
information for an arbitrary causal index and proves that the net generated by the Weyl
operators of localized sources satisfies generation, causality and time slice in every
Weyl system (isotony holds for every generated net), and covariance in the CCR algebra, with the composition law of Axiom~5
of~\cite{dimock}.}}

\subsection{Phase-space data}

\begin{definition}[Phase-space data]\label{def:phase-space}
Fix a causal index on a preordered set $\cR$ of regions. \emph{Phase-space data} consist of real
vector spaces $\cD$ (test functions) and $V$ (the phase space), a monotone map
$O\mapsto\cD(O)$ from regions to linear subspaces of $\cD$ such that every $f\in\cD$ lies in
some $\cD(O)$, an alternating bilinear form $\sigma$ on $V$, and a surjective linear map
$f\mapsto[f]$ from $\cD$ onto $V$, which assigns to each test function the element of the phase
space it generates, such that:
\begin{enumerate}[label=(P\arabic*),leftmargin=3em]
\item if $O_1\perp O_2$, $f\in\cD(O_1)$ and $g\in\cD(O_2)$, then $\sigma([f],[g])=0$;
\item if $O\dep O'$ and $f\in\cD(O)$, there is $f'\in\cD(O')$ with $[f']=[f]$.
\end{enumerate}
\end{definition}
\noindent\motiv{\vet{``Phase-space data'' is a term of these notes; the object of the
literature is a symplectic space of solutions together with the test functions localized in each
region~\cite{bgp}. We require $f\mapsto[f]$ only to be surjective, so $V$ may be a
quotient of $\cD$ or a space of solutions, and $\sigma$ may be degenerate: nothing in this
\revh{chapter} uses more, and \revf{Chapter~\ref{sec:simple} adds nondegeneracy
(Definition~\ref{def:nondegenerate}) as a hypothesis}.} \revc{For the free field,
$V$ with $\sigma$ is the covariant phase space of physics, and $\sigma([f],[g])=\langle
f,Eg\rangle$ is the Peierls bracket of the smeared fields $\phi(f)$ and
$\phi(g)$~\cite{peierls}, whose quantization is $[\phi(f),\phi(g)]=\ii\langle
f,Eg\rangle$ (Remark~\ref{rem:dictionary}).}}

\begin{extra}
\begin{example}[The lattice field]\label{ex:lattice-phase}
On the lattice causal index of Example~\ref{ex:lattice-index} take $\cD$ the finitely supported
$f:\Z^2\to\R$, $\cD(O)=\{f:\supp f\subseteq O\}$, $V=\cD/\ker E$ with $E$ the \revd{causal propagator} of
Definition~\ref{def:green}, $[f]$ the class of $f$, and $\sigma([f],[g])=\langle f,Eg\rangle$.
Every $f$ lies in $\cD(\supp f)$, since every subset of $\Z^2$ is a region.
The form is well defined and alternating by Proposition~\ref{prop:commutator}(b):
$\langle f,Eg\rangle=-\langle g,Ef\rangle$ vanishes when $Ef=0$ or $Eg=0$, and
$\langle f,Ef\rangle=0$. Condition (P1) is Proposition~\ref{prop:commutator}(c), since
orthogonality is antitone in both arguments, and (P2) is Corollary~\ref{cor:lattice-dep}, since
$Ef'=Ef$ means $[f']=[f]$. The form is nondegenerate: \revc{$\sigma([f],[\delta_p])=-\langle\delta_p,Ef\rangle=-(Ef)(p)$, so}
$\sigma([f],[\delta_p])=0$ for all $p$ gives $Ef=0$.
\end{example}
\end{extra}

\begin{definition}[Isomorphism of phase-space data]\label{def:phase-iso}
An \emph{isomorphism} $\theta=(\kappa,\kappa_*,L)$ from phase-space data
$(\cD,V,\sigma,[\cdot])$ to $(\cD',V',\sigma',[\cdot]')$ consists of an isomorphism $\kappa$ of
the causal indices (Definition~\ref{def:iso}), a linear bijection $\kappa_*:\cD\to\cD'$ and a
linear bijection $L:V\to V'$ with
\[
\kappa_*(\cD(O))=\cD'(\kappa O),\qquad [\kappa_*f]'=L[f],\qquad \sigma'(Lu,Lv)=\sigma(u,v).
\]
\end{definition}
\noindent\motiv{\vet{The last condition makes $L$ symplectic, as in \revh{Chapter}~\ref{sec:weyl}. It is
where the time orientation enters: a time reversal exchanges the retarded and advanced Green
operators (Remark~\ref{rem:time-orientation}), so the triple formed by it, the pushforward of
test functions and the induced map $L$ has $\sigma'(Lu,Lv)=-\sigma(u,v)$ and is not an isomorphism
when $\sigma\neq0$.}\extraonly{ \revc{Other linear maps over a time reversal are not excluded:
Remark~\ref{rem:massless-reversal} gives one that is an isomorphism.}}}

\begin{extra}
\begin{remark}[Time reversal of the massless lattice field]\label{rem:massless-reversal}
\revc{Let $m^2=0$ in Example~\ref{ex:lattice-phase}, let $r(t,x)=(-t,x)$ be the time reversal of
Example~\ref{ex:lattice-aut}, let $\varepsilon(t,x)=(-1)^{t+x}$, and put
$\kappa_*f=\varepsilon\cdot(f\circ r)$. Then $\kappa_*(\cD(O))=\cD(rO)$, and $\kappa_*\kappa_*f=f$.
Since $G^-$ is defined by the time-reflected kernel, $G^\pm(f\circ r)=(G^\mp f)\circ r$, hence
$E(f\circ r)=-(Ef)\circ r$. By Proposition~\ref{prop:commutator}(d), $\Delta(p-q)\neq0$ implies
$\varepsilon(p)\varepsilon(q)=-1$, hence $E(\varepsilon g)=-\varepsilon\,Eg$. Together,
$E\kappa_*f=\varepsilon\cdot((Ef)\circ r)$. So $\kappa_*$ maps $\ker E$ onto itself,
$L[f]=[\kappa_*f]$ is a well-defined linear bijection of $V$, and
\[
\sigma(L[f],L[g])=\langle\varepsilon\cdot(f\circ r),\varepsilon\cdot((Eg)\circ r)\rangle
=\langle f,Eg\rangle=\sigma([f],[g]).
\]
Thus $\theta=(r,\kappa_*,L)$ is an isomorphism of phase-space data with $\theta\theta=\mathrm{id}$,
and by Theorem~\ref{thm:ccr-cov} $(r,\alpha_\theta)$ is a covariance pair of the CCR net of the
massless lattice field with $\alpha_\theta^2=\mathrm{id}$, so the net is covariant under the group
$\{\mathrm{id},r\}$. For $m^2\neq0$ the
sign $\varepsilon$ no longer anticommutes with $E$ \revd{(the recursion gives $K(2,0)=-m^2$, so $\Delta(2,0)=m^2$ at a point with $t+x$ even, and
$E(\varepsilon\delta_{(0,0)})(2,0)=m^2\neq-m^2=-\varepsilon(2,0)\,(E\delta_{(0,0)})(2,0)$)}, and these notes do not decide whether another
isomorphism over the time reversal exists.}
\end{remark}
\end{extra}

\begin{lemma}[Identity and composite]\label{lem:phase-iso}
\revp{For phase-space data $(\cD,V,\sigma,[\cdot])$ the identity maps form an isomorphism from
these data to themselves. Let $\theta_1=(\kappa_1,\kappa_{1*},L_1)$ be an isomorphism from
$(\cD,V,\sigma,[\cdot])$ to $(\cD',V',\sigma',[\cdot]')$ and $\theta_2=(\kappa_2,\kappa_{2*},L_2)$
an isomorphism from $(\cD',V',\sigma',[\cdot]')$ to $(\cD'',V'',\sigma'',[\cdot]'')$. Then
$\theta_2\theta_1=(\kappa_2\kappa_1,\kappa_{2*}\kappa_{1*},L_2L_1)$ is an isomorphism from
$(\cD,V,\sigma,[\cdot])$ to $(\cD'',V'',\sigma'',[\cdot]'')$.}
\end{lemma}
\begin{proof}
$\kappa_{2*}\kappa_{1*}(\cD(O))=\kappa_{2*}(\cD'(\kappa_1O))=\cD''(\kappa_2\kappa_1O)$;
$[\kappa_{2*}\kappa_{1*}f]''=L_2[\kappa_{1*}f]'=L_2L_1[f]$; and
$\sigma''(L_2L_1u,L_2L_1v)=\sigma'(L_1u,L_1v)=\sigma(u,v)$.
\end{proof}

\subsection{The net of a Weyl system}

\noindent\vet{From here on we fix phase-space data $(\cD,V,\sigma,[\cdot])$ on a causal index;
primed data $(\cD',V',\sigma',[\cdot]')$ denote a second such choice\revc{, possibly on another
causal index}.}

\begin{definition}[The net of a Weyl system]\label{def:weyl-net}
Let $(A,W)$ be a Weyl system over $(V,\sigma)$ (Definition~\ref{def:weyl-system}). Its generators
are $G(O)=\{W([f]):f\in\cD(O)\}$, a family of generators in the sense of
Definition~\ref{def:generated}, and its \emph{net} is the net they generate,
$\cA_W(O)=\calg\,G(O)$.
\end{definition}

\begin{theorem}[Generation, causality, time slice]\label{thm:weyl-net}
For \revk{phase-space data $(\cD,V,\sigma,[\cdot])$ on a causal index and} every Weyl system $(A,W)$
over $(V,\sigma)$\revk{,} the net $\cA_W$ is generating, causal and has the time-slice property.
\end{theorem}
\begin{proof}
The three properties follow from Proposition~\ref{prop:generated} once we check the
generators. Generation: $\bigcup_OG(O)=W(V)$, since every $u\in V$ is some $[f]$ and every $f$ lies
in some $\cD(O)$; and $W(V)$ generates $A$. Causality: for $O_1\perp O_2$, $f\in\cD(O_1)$,
$g\in\cD(O_2)$, (P1) gives $\sigma([f],[g])=0$, so $W[f]$ and $W[g]$ commute by
Lemma~\ref{lem:weyl-system}; the adjoint $W[g]^*=W[-g]$ is again a generator of $O_2$, so the
same argument applies to it. Time slice: for $O\dep O'$ and $f\in\cD(O)$, (P2) gives
$f'\in\cD(O')$ with $W[f]=W[f']\in G(O')$.
\end{proof}

\begin{extra}
\noindent\motiv{\vet{For the lattice field of Example~\ref{ex:lattice-phase} and the CCR algebra
this \revm{gives generation, causality and time slice} without
assumptions on the spacetime; \reve{Chapter~\ref{sec:latticenet} completes that instance with
translation covariance}.}}
\end{extra}

\subsection{Covariance}

\begin{theorem}[Covariance pairs from Weyl systems]\label{thm:weyl-cov}
Let $\theta=(\kappa,\kappa_*,L)$ be an isomorphism of phase-space data, $(A,W)$ and $(A',W')$ Weyl
systems over $(V,\sigma)$ and $(V',\sigma')$, and $\alpha:A\to A'$ a $*$-isomorphism with
$\alpha(W(u))=W'(Lu)$ for all $u$. Then $(\kappa,\alpha)$ is a covariance pair from $\cA_W$ to
$\cA'_{W'}$.
\end{theorem}
\begin{proof}
By Proposition~\ref{prop:generated}(d) it suffices that $\alpha(G(O))=G'(\kappa O)$, and
$\alpha(G(O))=\{W'(L[f]):f\in\cD(O)\}=\{W'([\kappa_*f]'):f\in\cD(O)\}=\{W'([f']'):f'\in
\cD'(\kappa O)\}$.
\end{proof}
\noindent\motiv{\vet{Only the first two conditions of Definition~\ref{def:phase-iso} enter
here, since $\alpha$ is given; \revc{the symplectic condition is preserved under composition
(Lemma~\ref{lem:phase-iso}) and is used in Theorem~\ref{thm:ccr-cov}, where it makes
$\alpha_\theta$ exist. When $A'\neq0$ it also follows from the hypotheses: applying $\alpha$ to
the Weyl relation gives $e^{-\ii\sigma(u,w)/2}=e^{-\ii\sigma'(Lu,Lw)/2}$, and the scaling argument
of Remark~\ref{rem:antisymplectic} gives $\sigma'(Lu,Lw)=\sigma(u,w)$.}}}

\begin{corollary}[Isomorphic Weyl systems]\label{cor:weyl-iso}
If $(A,W)$ and $(A',W')$ are Weyl systems over $(V,\sigma)$, both nets formed from the same
phase-space data, and $\alpha:A\to A'$ is a
$*$-isomorphism with $\alpha\circ W=W'$, the nets $\cA_W$ and $\cA'_{W'}$ are isomorphic.
\end{corollary}
\begin{proof}
Theorem~\ref{thm:weyl-cov} with the identity isomorphism.
\end{proof}
\noindent\motiv{\vet{For nondegenerate $\sigma$ and nonzero algebras such an $\alpha$ always
exists, by the uniqueness theorem for the CCR algebra~\cite{slawny} (see also~\cite{manuceau}); the paper
\cite[p.~225]{dimock} uses it, by citation, for the independence of the net from the
representation in its Theorem~3. These notes do not prove it and do not formalize it here; the
corollary carries $\alpha$ as a hypothesis.}}

\begin{definition}[The CCR net]\label{def:ccr-net}
The \emph{CCR net} of phase-space data is $\cA_W$ for the Weyl system $(\CCR(V,\sigma),W)$ of
Proposition~\ref{prop:ccr-system}. For an isomorphism $\theta=(\kappa,\kappa_*,L)$ let
$\alpha_\theta:\CCR(V,\sigma)\to\CCR(V',\sigma')$ be conjugation by the relabelling unitary $U_L$
(Proposition~\ref{prop:relabel}).
\end{definition}

\begin{theorem}[Covariance of the CCR net]\label{thm:ccr-cov}
For every isomorphism $\theta\revk{=(\kappa,\kappa_*,L)}$ \revk{from} phase-space data
\revk{$(\cD,V,\sigma,[\cdot])$ to $(\cD',V',\sigma',[\cdot]')$}, $(\kappa,\alpha_\theta)$ is a covariance
pair \revk{from the CCR net of the first to that of the second}, $\alpha_{\mathrm{id}}=\mathrm{id}$, and
$\alpha_{\theta_2\theta_1}=\alpha_{\theta_2}\alpha_{\theta_1}$.
\end{theorem}
\begin{proof}
Proposition~\ref{prop:relabel} gives $\alpha_\theta(W(u))=W'(Lu)$, $\alpha_{\mathrm{id}}=\mathrm{id}$
and $\alpha_{L_2L_1}=\alpha_{L_2}\alpha_{L_1}$, and the symplectic map of $\theta_2\theta_1$ is
$L_2L_1$ (Lemma~\ref{lem:phase-iso}). Theorem~\ref{thm:weyl-cov} gives the covariance pair.
\end{proof}
\noindent\motiv{\vet{This is Axiom~5 of~\cite[p.~220]{dimock} for the CCR net: an assignment
$\theta\mapsto\alpha_\theta$ with $\alpha_{\mathrm{id}}=\mathrm{id}$ and
$\alpha_{\theta_2\theta_1}=\alpha_{\theta_2}\alpha_{\theta_1}$, each $(\kappa,\alpha_\theta)$ a
covariance pair (Remark~\ref{rem:axiom5}), indexed by the isomorphisms of phase-space data. \revc{We obtain an assignment indexed by the isometries
themselves, as in~\cite{dimock}, by taking for each isometry the triple of the
induced map of regions, the pushforward of test functions and the induced map on $V$; pushforwards
compose, and $L$ is determined by $\kappa_*$ because $[\cdot]$ is surjective, so this choice
respects composition.}}}

\motiv{\vet{For the
free field these include the maps induced by the isometries that preserve the time orientation,
which commute with the Green operators by their uniqueness on a globally hyperbolic
spacetime~\cite[Cor.~3.4.3]{bgp}, a result \revk{formalized here only for the massless field on
Minkowski space (Theorem~\ref{thm:mink-unique}, Proposition~\ref{prop:mink-symmetry})} \revlt{and for
the Klein--Gordon field there (Theorem~\ref{thm:kg-field}(b), (c))}; the converse fails, since
$(\mathrm{id},-\mathrm{id},-\mathrm{id})$ is an isomorphism that no isometry induces \revc{when $\cD\neq0$}. Covariance under a group of
such isomorphisms is \revk{\extraonly{Theorem~\ref{thm:lattice-net} for the lattice translations and
}Theorem~\ref{thm:mink-covariance} \revlt{and Theorem~\ref{thm:kg-field}(c)} for the orthochronous Poincar\'e group}.}}
\endgroup
